\begin{figure}[t]
\centering
% Round timelines (schematic, not to scale).
\begin{tikzpicture}[font=\scriptsize,x=0.74cm,y=0.36cm]
\definecolor{freshc}{RGB}{35,139,69}
\definecolor{lookc}{RGB}{199,233,192}
\definecolor{verc}{RGB}{49,130,189}
% \blk{x0}{x1}{y}{label}{fill}{text colour}
\newcommand{\blk}[6]{%
  \filldraw[fill=#5,draw=black!70,line width=0.3pt] (#1,#3) rectangle (#2,{#3+0.8});
  \node[text=#6] at ({(#1+#2)/2},{#3+0.4}) {#4};}
% \miss{x0}{x1}{y}: discarded lookahead
\newcommand{\miss}[3]{\draw[red!75!black,line width=0.6pt] (#1,#3) -- (#2,{#3+0.8}) (#1,{#3+0.8}) -- (#2,#3);}
% \rows{y of edge row}{lane title}
\newcommand{\rows}[2]{%
  \node[anchor=east] at (-0.05,{#1+0.4}) {Edge};
  \node[anchor=east] at (-0.05,{#1-0.6}) {Cloud};
  \node[anchor=west,font=\scriptsize\itshape] at (-1.2,{#1+1.8}) {#2};
  \draw[black!35] (0,{#1-1.15}) -- (9.6,{#1-1.15});}
% (a) synchronous
\rows{9}{(a) Sync-4: the device idles while the cloud verifies}
\blk{0}{1.6}{9}{draft 4}{freshc}{white}
\blk{1.6}{3.6}{8}{net + verify}{verc}{white}
\blk{3.6}{5.2}{9}{draft 4}{freshc}{white}
\blk{5.2}{7.2}{8}{net + verify}{verc}{white}
\blk{7.2}{8.8}{9}{draft 4}{freshc}{white}
% (b) fixed lookahead
\rows{5}{(b) Fixed-4: a miss redrafts the whole window}
\blk{0}{1.6}{5}{draft 4}{freshc}{white}
\blk{1.6}{3.6}{4}{verify}{verc}{white}
\blk{1.6}{3.6}{5}{lookahead 5}{lookc}{black}
\blk{3.6}{5.6}{4}{verify}{verc}{white}
\blk{3.6}{5.6}{5}{lookahead 5}{lookc}{black}
\miss{3.6}{5.6}{5}
\blk{5.6}{7.2}{5}{draft 4}{freshc}{white}
\blk{7.2}{9.2}{4}{verify}{verc}{white}
\node[font=\tiny] at (3.6,6.15) {hit};
\node[font=\tiny] at (5.6,6.15) {miss};
% (c) CAPE
\rows{1}{(c) CAPE: short fresh window, lookahead fills the round trip}
\blk{0}{0.6}{1}{$\gamma_f$}{freshc}{white}
\blk{0.6}{2.6}{0}{verify}{verc}{white}
\blk{0.6}{2.6}{1}{lookahead $h$}{lookc}{black}
\blk{2.6}{4.6}{0}{verify}{verc}{white}
\blk{2.6}{4.6}{1}{lookahead $h$}{lookc}{black}
\miss{2.6}{4.6}{1}
\blk{4.6}{5.2}{1}{$\gamma_f$}{freshc}{white}
\blk{5.2}{7.2}{0}{verify}{verc}{white}
\blk{5.2}{7.2}{1}{lookahead $h$}{lookc}{black}
\blk{7.2}{9.2}{0}{verify}{verc}{white}
\node[font=\tiny] at (2.6,2.15) {hit};
\node[font=\tiny] at (4.6,2.15) {miss};
\draw[-{Latex[length=1.6mm]},black!60] (0,-0.55) -- (9.6,-0.55) node[below left,font=\tiny] {time};
\end{tikzpicture}

\caption{Round timelines (schematic). Fresh drafting (dark) delays the next verification; lookahead (light) overlaps it and is reused after a hit or discarded after a miss (crossed). \cape{} drafts few fresh tokens and sizes the lookahead to the round trip.}
\label{fig:timeline}
\end{figure}

\section{Background}
\label{sec:measure}

\subsection{Edge--cloud speculative decoding}
Let $x$ be the current prefix. In each round the device drafts $\gamma$ tokens $D=(d_1,\dots,d_\gamma)$ with the draft model and sends them to the cloud. The target model scores $x\Vert D$ in one forward pass and returns the number $n$ of leading draft tokens that equal its own greedy predictions, together with one target token $t$. This token is the correction at the first mismatch, or a bonus token if all $\gamma$ draft tokens are accepted. The prefix becomes $x\Vert d_{1:n}\Vert t$, so each round emits $n+1$ tokens. We use strict greedy acceptance: a draft token is accepted only if it equals the target's highest-scoring token, so every emitted token is the target model's greedy choice.

\emph{Lookahead.} While the cloud verifies $D$, the device can draft $h$ further tokens $L=(\ell_1,\dots,\ell_h)$ from $x\Vert D$, assuming full acceptance; $\ell_1$ is the draft model's guess of the bonus token. If $n=\gamma$ and $t=\ell_1$, then $\ell_{2:h}$ is a valid draft for the next round and can be sent at once (a \emph{hit}). Otherwise $L$ is discarded (a \emph{miss}). A fixed window with $h=\gamma+1$ prefetches exactly one complete next block (Fig.~\ref{fig:timeline}(b)).

\emph{Interruption.} If the verdict rejects one of the sent tokens, the lookahead drafted from them is useless before it is finished. PicoSpec interrupts it as soon as the verdict arrives~\cite{picospec2026}, which frees the draft server for the next round. Every pipelined fixed window in our evaluation, and \cape{}, interrupts in this way.

\subsection{Testbed}
\label{sec:testbed}
\emph{Edge device.} The edge device is an NVIDIA Jetson Orin Nano Super developer kit (8\,GB) in its 25\,W power mode, running JetPack 6.2.1 (L4T R36.4.7). It serves the draft model with llama.cpp (build b9744) as a single-slot HTTP server with all layers on the GPU. Board input power is sampled at 10\,Hz from the INA3221 VDD\_IN rail and integrated over each request. This is the gross energy the device pays, including idle and operating-system power.

\emph{Cloud.} The cloud verifier runs on a server with an AMD Ryzen Threadripper PRO 5965WX CPU (24 cores), 125\,GiB of memory and an NVIDIA GeForce RTX 3080 GPU (10\,GB), running Ubuntu 24.04 with NVIDIA driver 580.173.02, vLLM 0.26.0~\cite{vllm2023} and PyTorch 2.11.0. A token-ID verification API in front of vLLM returns the target's top-1 token at every draft position, with prefix caching enabled; it is identical for all methods. The server runs no other workload.

\emph{Model pairs.} Table~\ref{tab:pairs} lists the four draft/target pairs. They span three model families and draft models from 270M to 2B parameters.

\begin{table}[t]
\centering
\caption{Draft/target pairs. Draft models run on the Jetson with llama.cpp; targets run on the cloud GPU with vLLM.}
\label{tab:pairs}
\footnotesize
\setlength{\tabcolsep}{3pt}
\begin{tabular}{@{}llll@{}}
\toprule
Pair & Draft (Jetson) & Target (cloud) & Context\\
\midrule
Llama & Llama-3.2-1B-Instruct, Q4\_K\_M & Llama-3.1-8B-Instruct, AWQ & 4,096\\
Qwen & Qwen2.5-0.5B-Instruct, Q4\_K\_M & Qwen2.5-7B-Instruct, AWQ & 4,096\\
Gemma-2 & Gemma-2-2B-it, Q4\_K\_M & Gemma-2-9B-it, AWQ & 2,048\\
Gemma-3 & Gemma-3-270M-it, Q8\_0 & Gemma-3-4B-it, FP8 & 2,048\\
\bottomrule
\end{tabular}
\end{table}

\emph{Network.} Both the Jetson and the server are connected to the campus Wi-Fi network, so every verification crosses two wireless links: a fast path that carries the jitter of both. To emulate wide-area round trips, the client injects an additional delay into every verification request on top of this real path (Section~\ref{sec:eval}).

\subsection{Where the cost of speculation falls}
Three properties of the edge--cloud loop shape our design. \emph{Fresh drafting is on the critical path:} the verifier cannot start before the last draft token of a round has arrived, so every token drafted after a verification returns delays the next one. \emph{Verification is priced per round:} a GPU target scores all draft positions of a round in one forward pass, so the cost of a draft token is paid on the device. \emph{Draft work has an energy price:} the board's energy per request is
\begin{equation}
E \approx P_0\, T + P_d\, B, \label{eq:energy}
\end{equation}
where $T$ is the request's wall time, $B$ the draft server's busy time, and $P_0$ and $P_d$ the base and draft-active power. Lookahead drafted during verification is free in $T$ but not in $B$, and a discarded block raises the energy bill without emitting a token. Interruption bounds the time that a doomed lookahead costs by the round trip, but its energy is spent until the cut.

These properties lead to four goals. (i)~Keep fresh drafting off the critical path. (ii)~Fill the verification time with lookahead, but only as much as is likely to be reused, and less when energy matters. (iii)~Adapt to the round trip, its jitter and the acceptance rate, which vary across and within requests. (iv)~Leave the cloud server and the verification rule unchanged.
